\documentclass[11pt,a4paper]{book}

% -------------------- Packages --------------------
\usepackage[utf8]{inputenc}
\usepackage{amsmath, amssymb, amsthm}
\usepackage{graphicx}
\usepackage{geometry}
\usepackage{hyperref}
\usepackage{listings}
\usepackage{xcolor}
\usepackage{tcolorbox}
\usepackage{fancyhdr}

\geometry{margin=2.5cm}

% -------------------- Style --------------------
\definecolor{mainblue}{RGB}{30,60,120}
\definecolor{lightgray}{RGB}{245,245,245}

\hypersetup{
    colorlinks=true,
    linkcolor=mainblue,
    urlcolor=mainblue
}

\pagestyle{fancy}
\fancyhf{}
\fancyhead[L]{Algorithms I}
\fancyhead[R]{Lecture Notes}
\fancyfoot[C]{\thepage}

% -------------------- Boxes --------------------
\newtcolorbox{definitionbox}{colback=lightgray,colframe=mainblue,title=Definition}
\newtcolorbox{theorembox}{colback=white,colframe=mainblue,title=Theorem}
\newtcolorbox{warningbox}{colback=white,colframe=red!70!black,title=Warning}

% -------------------- Code --------------------
\lstset{
    basicstyle=\ttfamily\small,
    keywordstyle=\color{mainblue},
    commentstyle=\color{gray},
    frame=single,
    breaklines=true
}

% -------------------- Document --------------------
\begin{document}

\chapter{Algorithms Fundamentals}

\section{Definition of an Algorithm}
\begin{definitionbox}
An algorithm is a well-defined computational procedure that takes input values and produces output values through a sequence of computational steps.
\end{definitionbox}

\section{Asymptotic Complexity}

\begin{itemize}
    \item Big-O: upper bound
    \item Big-Omega: lower bound
    \item Big-Theta: tight bound
    \item little-o: strict upper bound
    \item little-omega: strict lower bound
\end{itemize}

\section{Sorting Algorithms}

\subsection{Insertion Sort}
\textbf{Complexity:} $O(n^2)$ \\n\textbf{Space:} $O(1)$

\begin{lstlisting}
Insertion-Sort(A, n)
    for j = 2 to n
        key = A[j]
        i = j - 1
        while i > 0 and A[i] > key
            A[i+1] = A[i]
            i = i - 1
        A[i+1] = key
\end{lstlisting}

\subsection{Merge Sort}
\textbf{Complexity:} $O(n \log n)$ \\n\textbf{Space:} $O(n)$

\begin{lstlisting}
MERGE-SORT(A,p,r)
    if p < r
        q = (p+r)/2
        MERGE-SORT(A,p,q)
        MERGE-SORT(A,q+1,r)
        MERGE(A,p,q,r)
\end{lstlisting}

\section{Divide and Conquer}

\subsection{Maximum Subarray}
\textbf{Complexity:} $O(n \log n)$

\begin{lstlisting}
FIND-MAXIMUM-SUBARRAY(A, low, high)
    if high == low
        return A[low]
    mid = (low + high)/2
    left = ...
    right = ...
    cross = ...
    return max(left,right,cross)
\end{lstlisting}

\section{Matrix Multiplication}

\subsection{Naive Recursive Method}
\textbf{Complexity:} $O(n^3)$

\begin{lstlisting}
REC-MAT-MULT(A,B,n)
    if n == 1
        return A*B
    split matrices
    compute subproblems
\end{lstlisting}

\subsection{Strassen Algorithm}
\textbf{Complexity:} $O(n^{\log_2 7})$

\section{Karatsuba Multiplication}
\textbf{Complexity:} $O(n^{\log_2 3})$

Split numbers:
\[
x = x_H b^{n/2} + x_L, \quad y = y_H b^{n/2} + y_L
\]

\section{Heaps}

\subsection{Heapify}
\textbf{Complexity:} $O(\log n)$

\begin{lstlisting}
MAX-HEAPIFY(A,i,n)
    l = LEFT(i)
    r = RIGHT(i)
    ...
\end{lstlisting}

\subsection{Heapsort}
\textbf{Complexity:} $O(n \log n)$

\section{Priority Queue}
Operations:
\begin{itemize}
    \item Insert: $O(\log n)$
    \item Extract-Max: $O(\log n)$
    \item Maximum: $O(1)$
\end{itemize}

\section{Binary Search Trees}

\textbf{Search Complexity:} $O(h)$

\begin{lstlisting}
TREE-SEARCH(x,k)
    if x == NIL or k == x.key
        return x
    if k < x.key
        return left
    else right
\end{lstlisting}

\section{Tree Operations}
\begin{itemize}
    \item Insert: $O(h)$
    \item Delete: $O(h)$
    \item Successor: $O(h)$
\end{itemize}

\section{Data Structures Comparison}

\begin{itemize}
    \item Stack: LIFO, $O(1)$
    \item Queue: FIFO, $O(1)$
    \item Linked List: $O(n)$ search
    \item BST: $O(h)$ operations
\end{itemize}

\chapter{Le quiz d'aujourd'hui}
Efficient algorithms depend heavily on choosing the correct data structure.

\end{document}
